\begin{figure}[htbp]
\centering
\begin{tikzpicture}[font=\footnotesize]

% ---- forecasts under comparison ----
\node[fgout=3.2cm] (own) at (-3.2, 0)
  {\textbf{OBTF forecasts}\\\(\hat{L}\), \(\hat{S}\), \(\hat{W}\)};
\node[fgdata=4.5cm, aspect=0.07] (ent) at (3.2, 0)
  {\textbf{\mbox{ENTSO-E} forecasts}\\\mbox{\(L^{\mathrm{ENTSOE}}\), \(S^{\mathrm{ENTSOE}}\), \(W^{\mathrm{ENTSOE}}\)}};

% ---- evaluation row ----
\node[fgeval=3.6cm] (eval) at (0,-2.6)
  {\textbf{Evaluation}\\RMSE, MAE, one-sided\\Giacomini--White test};
\node[fgdata=2.2cm] (real) at (4.4,-2.6)
  {\textbf{Realized values}\\\(L\), \(S\), \(W\)};

\node[fgrq=6.4cm] (rq) at (0,-4.5)
  {\textbf{RQ1:} How do the OBTF forecasts compare with the
   \mbox{ENTSO-E} forecasts?};

% ---- flows ----
\draw[fgflow] (own.south) -- ++(0,-0.55) -| ([xshift=-0.9cm]eval.north);
\draw[fgflow] (ent.south) -- ++(0,-0.55) -| ([xshift=0.9cm]eval.north);
\draw[fgflow] (real.west) -- (eval.east);
\draw[fgflow] (eval.south) -- (rq.north);

\end{tikzpicture}
\caption[OBTF forecast comparison for RQ1.]
{Design of RQ1. The OBTF load, solar, and wind
forecasts and the \mbox{ENTSO-E} forecasts are each scored against the realized
values on a matched panel. The OBTF forecasts use only information
available before the 12:00 deadline. Notation as in
Figure~\ref{fig:system_overview}.}
\label{fig:rq1_component_forecast_design}
\end{figure}
